%% Đề 21 -- de-thi-thu-tuyen-sinh-lop-10-mon-toan-lan-1-nam-2026-so-gddt-lang-son.pdf
\deheader{ĐỀ SỐ 21}%
  {Kỳ thi Tuyển sinh vào lớp 10 THPT — Năm 2026}%
  {Sở GD\&ĐT Lạng Sơn | Chỉ phần Đại số}
\addcontentsline{toc}{section}{Đề 21 — Tuyển sinh vào lớp 10 THPT — Năm 2026}

\begin{de_cau}{Câu 1}{2,0 điểm}
\begin{enumerate}[label=\alph*)]
\item Tính: $A = \sqrt{48} + \sqrt{75} - 2\sqrt{108}$.
\item Cho $P = \dfrac{x+2}{x\sqrt{x}+1} + \dfrac{\sqrt{x}-1}{x-\sqrt{x}+1} - \dfrac{2}{\sqrt{x}+1}$, $x\geq 0$.
Rút gọn $P$. Tính $P$ khi $x = 4$.
\end{enumerate}
\end{de_cau}

\begin{loi_giai_cau}
\textbf{a)} $4\sqrt{3}+5\sqrt{3}-12\sqrt{3}=(4+5-12)\sqrt{3}=-3\sqrt{3}$.

\textbf{b)} $x\sqrt{x}+1=(\sqrt{x}+1)(x-\sqrt{x}+1)$.
\[P=\frac{x+2+(\sqrt{x}-1)(\sqrt{x}+1)-2(x-\sqrt{x}+1)}{(\sqrt{x}+1)(x-\sqrt{x}+1)}\]
\[=\frac{x+2+x-1-2x+2\sqrt{x}-2}{(\sqrt{x}+1)(x-\sqrt{x}+1)}=\frac{2\sqrt{x}-1}{(\sqrt{x}+1)(x-\sqrt{x}+1)}.\]

Khi $x=4$: $P=\dfrac{2\cdot 2-1}{3\cdot 3}=\dfrac{3}{9}=\dfrac{1}{3}$.
\end{loi_giai_cau}

\begin{de_cau}{Câu 2}{2,0 điểm}
\begin{enumerate}[label=\alph*)]
\item Giải hệ phương trình: 
\[ \begin{cases}3x-2y=7\\x+y=3\end{cases} \].
\item Cho hàm số bậc hai $y=ax^2$ ($a\neq 0$) đi qua điểm $B(2;-2)$. Tìm $a$ và vẽ đồ thị.
\end{enumerate}
\end{de_cau}

\begin{loi_giai_cau}
\textbf{a)} Từ (2): $y=3-x$. Thế (1): $3x-2(3-x)=7\Rightarrow 5x=13\Rightarrow x=\frac{13}{5}, y=\frac{2}{5}$.

\textbf{b)} $-2=a\cdot 4\Rightarrow a=-\frac{1}{2}$. Hàm số $y=-\frac{1}{2}x^2$, parabol đỉnh $O$, lõm xuống.
\end{loi_giai_cau}

\begin{de_cau}{Câu 3}{2,0 điểm}
Quãng đường từ A đến B dài 180 km. Một ô tô đi từ A đến B với vận tốc $v$ km/h. Do mưa, xe phải
giảm vận tốc 10 km/h trên 60 km cuối, nên đến B chậm hơn dự định 24 phút. Tính $v$.
\end{de_cau}

\begin{loi_giai_cau}
Thời gian dự định: $\dfrac{180}{v}$ giờ.

Thời gian thực: $\dfrac{120}{v}+\dfrac{60}{v-10}$ giờ.

Phương trình:
\[\frac{120}{v}+\frac{60}{v-10}-\frac{180}{v}=\frac{24}{60}=\frac{2}{5}.\]
\[\frac{60}{v-10}-\frac{60}{v}=\frac{2}{5}\Rightarrow\frac{600}{v(v-10)}=\frac{2}{5}\Rightarrow v^2-10v-1500=0.\]
$(v-\frac{10\pm\sqrt{100+6000}}{2})$: $v=\dfrac{10+\sqrt{6100}}{2}\approx 44$ km/h.

Nghiệm nguyên đẹp: $(v-40)(v+30)=0+\ldots$ Thử $v=40$: $40\cdot 30=1200\neq 1500$.
Thử lại: $v^2-10v=1500$, $v=45$: $2025-450=1575\neq 1500$. $v=40$: $1600-400=1200$.

Dùng công thức: $v=\dfrac{10+\sqrt{6100}}{2}\approx 44{,}1$ km/h.
\end{loi_giai_cau}

\begin{de_cau}{Câu 4}{1,5 điểm}
Cho phương trình $x^2-2(m+1)x+m^2+4m+3=0$ (tham số $m$).
\begin{enumerate}[label=\alph*)]
\item Chứng minh phương trình luôn có hai nghiệm phân biệt.
\item Tìm $m$ để $x_1^2+x_2^2-x_1x_2=7$.
\end{enumerate}
\end{de_cau}

\begin{loi_giai_cau}
\textbf{a)} $\Delta'=(m+1)^2-(m^2+4m+3)=m^2+2m+1-m^2-4m-3=-2m-2=-2(m+1)$.

Để $\Delta'>0$: $m<-1$. Vậy phương trình chỉ có 2 nghiệm phân biệt khi $m<-1$ (không luôn luôn).

Với $m=-1$: $\Delta'=0$, có nghiệm kép. Vậy phương trình có 2 nghiệm $\Leftrightarrow m<-1$.

\textbf{b)} Vi-ét: $S=2(m+1)$, $P=m^2+4m+3=(m+1)(m+3)$.

$x_1^2+x_2^2-x_1x_2=S^2-3P=4(m+1)^2-3(m+1)(m+3)=(m+1)[4(m+1)-3(m+3)]=(m+1)(m-5)=7$.

$m^2-4m-5=7\Rightarrow m^2-4m-12=0\Rightarrow(m-6)(m+2)=0$.

$m=6$ (loại, $m<-1$) hoặc $m=-2$ (thỏa). Vậy $\boxed{m=-2}$.
\end{loi_giai_cau}
